\documentclass[acmsmall,screen,nonacm]{acmart}
\usepackage{mathtools}
\usepackage{xr}
\externaldocument[P-]{../build/main}
\newcommand{\archivedoi}{10.5281/zenodo.24494014}
\newcommand{\wt}{\operatorname{wt}}
% Shared notation (moved from Section 3)
\newcommand{\NN}{\mathbb{N}}
\newcommand{\QQ}{\mathbb{Q}}
\newcommand{\RR}{\mathbb{R}}
\newcommand{\supp}{\operatorname{supp}}
\newcommand{\sgn}{\operatorname{sgn}}
\newcommand{\core}{\operatorname{core}}
\newcommand{\SC}{\operatorname{sc}}
\newcommand{\FQ}{\mathcal F_{\QQ}}
\newcommand{\FR}{\mathcal F_{\RR}}


\renewcommand{\thetable}{S\arabic{table}}
\title{Lean Correspondence Tables: Supplement to ``One Word Against Two: The Complexity of Revising Markov Laws Under Held Tests''}
\author{Ioannis Tsiokos}
\orcid{0009-0009-7659-5964}
\affiliation{%
  \institution{Automorph Inc.}
  \city{Athens}
  \country{Greece}}
\email{ioannis@automorph.io}
\renewcommand{\shortauthors}{Ioannis Tsiokos}
\begin{document}
\maketitle
\noindent These tables identify the Lean declarations and scope of each
listed check. The Lean~4 development is archived with the paper under the
DOI \texttt{\archivedoi}. Cross-references to the paper use its numbering.
Tables S1--S9 are numbered within this supplement. In each row, the third
column states what Lean checks and the fourth what is proved only on
paper.

% lint: lean-table begin
\DeclareUrlCommand\leanid{\urlstyle{tt}}
\newcommand{\lc}{\raggedright\let\\\tabularnewline\hspace{0pt}}
\newcommand{\leanrow}{\addlinespace[2pt]}
\newcommand{\leanhead}{%
  \toprule
  \lc Result & \lc Lean declarations, by file & \lc What Lean checks & \lc Proved only on paper\\
  \midrule}

\begin{table}[tbp]
\scriptsize\setlength{\tabcolsep}{3pt}
\caption{Lean correspondence for Section~\ref{P-sec:problem}. Declarations
are named relative to the namespace \texttt{MovingProjection}. In a cell,
each group gives a file and, where needed, a namespace:
``\texttt{F.lean}, \texttt{N}: \texttt{a}, \texttt{b}'' names
\texttt{N.a} and \texttt{N.b} in the file \texttt{F.lean}. Declarations marked
private are internal proof steps.
The same conventions hold in
Tables~\ref{tab:lean-core}--\ref{tab:lean-exactness}.}\label{tab:lean-problem}
\begin{tabular}{@{}p{0.14\textwidth}p{0.33\textwidth}p{0.27\textwidth}p{0.20\textwidth}@{}}
\leanhead
\lc Equation~\eqref{P-eq:word-joint}, Lemma~\ref{P-lem:compile}
& \lc \leanid{NativeWords.lean}: \leanid{nativeWordEndpoint_factor}.
\leanid{EventExperiments.lean}, \leanid{NativeExperiments}: \leanid{singleton_joint_iff}.
\leanid{ParityKernel.lean}, \leanid{ParityBasis}: \leanid{word_identity_iff}.
\leanid{SingleAdditiveLicence.lean}, \leanid{SingleAdditiveLicence}: \leanid{joint_words_iff}.
& \lc The factorization of word probabilities; (a)~as equal word probabilities
and equal projected word kernels, with kernels compared through eigenvalue
products; (b)~for nonempty words, the probability equation and one equation
per probe.
& \lc None.\\
\leanrow
\lc Definition~\ref{P-def:system}
& \lc \leanid{Spectral.lean}: \leanid{MonomialTest}, \leanid{MarkBudget}, \leanid{Legal}.
\leanid{Contraction.lean}: \leanid{Contracts}, \leanid{Reduces}, \leanid{IsCore}, \leanid{core}.
\leanid{Components.lean}: \leanid{signedCore}.
\leanid{Local.lean}: \leanid{FineConnected}.
\leanid{Experiments.lean}, \leanid{NativeExperiments}: \leanid{protocol}.
\leanid{Capabilities.lean}, \leanid{Protocol}: \leanid{fiber}.
\leanid{FixedProfile.lean}: \leanid{spectrumBudgets}.
& \lc The algebraic definitions of Definition~\ref{P-def:system}(a)--(d), and
the related fine-connectivity definition of
Section~\ref{P-sec:fine-connectivity}. For (e), \leanid{NativeExperiments.protocol}
and \leanid{Protocol.fiber} supply the held word and calibration rows and
the reversal rows, and \leanid{spectrumBudgets} supplies the per-label
budgets. The Lean relation \leanid{FineConnected} admits an empty chain even
at an unlawful point; the classification supplies feasibility of the
endpoints as separate hypotheses. The Lean model type assumes $\kappa<1$;
with a nonempty label set, the existence of a feasible parameter vector
also implies $\kappa\ge0$.
& \lc Computing the core (Theorem~\ref{P-thm:certificates}(a)).\\
\leanrow
\lc Theorem~\ref{P-thm:fixed-profile}
& \lc \leanid{FixedProfile.lean}: \leanid{finePhysical_iff_signedCore}.
\leanid{Experiments.lean}, \leanid{NativeExperiments}: \leanid{exact_six_iff}, \leanid{compiled_rate_tests}.
\leanid{ComplexityTheorem.lean}, \leanid{EventExperiments}: \leanid{singleton_criterion}.
& \lc The fixed-profile signed-core equivalence. The abstract form assumes
homogeneous rate rows and positive ratios. The native equal-length tests
with positive ratios supply those hypotheses.
& \lc None.\\
\leanrow
\lc Theorem~\ref{P-thm:boundary}(i)
& \lc \leanid{ComplexityTheorem.lean}, \leanid{EventExperiments}: \leanid{complexity_boundary}, \leanid{singleton_certificate}.
Rows for Theorems~\ref{P-thm:certificates} and~\ref{P-thm:p-hard}.
& \lc The signed-core criterion, soundness and existence of certificates on
presented inputs, and the circuit equivalence.
& \lc Membership in P, certificate size and checking time; the log-space
reduction and Goldschlager's theorem.\\
\leanrow
\lc Theorem~\ref{P-thm:boundary}(ii)
& \lc \leanid{ComplexityTheorem.lean}, \leanid{EventExperiments}: \leanid{complexity_boundary}.
Rows for Section~\ref{P-sec:hard}.
& \lc The NCL equivalence at the supplied tolerance and for fine reachability,
for each length of the word-against-pair test; distinct constructed endpoints for distinct input orientations.
& \lc Polynomial-time encoding, binary output size and the polynomial bound
on $1/\varepsilon_0$, together with Hearn--Demaine's hardness theorem.\\
\leanrow
\lc Corollary~\ref{P-cor:algebraic}, P-complete part
& \lc \leanid{AlgebraicBoundary.lean}, \leanid{MonotoneCircuitMarkov.Circuit}: \leanid{algebraic_fine_iff_value}, \leanid{algebraic_endpoint_legal}, \leanid{algebraic_bounded}.
\leanid{CoreCertificate.lean}, \leanid{CoreCertificate}: \leanid{checkPair_sound}.
\leanid{CertificateComplete.lean}, \leanid{CoreCertificate}: \leanid{checkPair_complete}, \leanid{checkCore_complete_short}.
& \lc The circuit system: feasible endpoints, the bounded promise and
equivalence to the circuit value; certificate soundness, existence and the
stage bound.
& \lc LP search with its round, dimension and binary-length bounds; certificate encoding;
running time; the log-space transduction.\\
\leanrow
\lc Corollary~\ref{P-cor:algebraic}, PSPACE part
& \lc \leanid{AlgebraicBoundary.lean}, \leanid{AlgebraicConstraintLogic}: \leanid{lawful_iff}, \leanid{endpoint_lawful}, \leanid{bounded}, \leanid{chain_of_flips}, \leanid{flips_of_chain}, \leanid{fine_iff_flips}.
& \lc Exact semantics of the charge binomials, budgets and trinomials;
feasible endpoints; boundedness; legal flips give sequences at every
tolerance, and a sequence at tolerance one gives legal flips. The forward
direction clips the toggles to $[0,1]$.
& \lc The polynomial-time encoding; Hearn and Demaine's theorem.\\
\bottomrule
\end{tabular}
\end{table}

\begin{table}[tbp]
\scriptsize\setlength{\tabcolsep}{3pt}
\caption{Lean correspondence for Section~\ref{P-sec:core}.}\label{tab:lean-core}
\begin{tabular}{@{}p{0.14\textwidth}p{0.33\textwidth}p{0.27\textwidth}p{0.20\textwidth}@{}}
\leanhead
\lc Lemma~\ref{P-lem:restriction}
& \lc (a) \leanid{Spectral.lean}: \leanid{rational_contraction}.
\leanid{Local.lean}: \leanid{fineConnected_scale}, \leanid{fineConnected_to_core}.
\leanid{RealComponents.lean}: \leanid{real_contract_path} (private).
(b) \leanid{Contraction.lean}: \leanid{Contracts.restrict}, \leanid{Reduces.restrict}.
(c) \leanid{Contraction.lean}: \leanid{exists_core}, \leanid{IsCore.unique}, \leanid{IsCore.of_reduces}, \leanid{core_eq_of_reduces}, \leanid{core_eq_iff_common_reduct}.
& \lc Part (a): feasible rational samples and fine connectivity, and a
feasible real monomial path. Parts (b) and (c): restriction and the
least-support conclusions.
& \lc None.\\
\leanrow
\lc Lemma~\ref{P-lem:irreducible}
& \lc \leanid{CoreOrder.lean}: \leanid{irreducibleSupport_iff_core_eq}, \leanid{core_greatest_irreducible}, \leanid{core_eq_union_irreducible}.
& \lc As stated.
& \lc None.\\
\leanrow
\lc Figure~\ref{P-fig:peeling} (the staged example)
& \lc \leanid{StagedContraction.lean}, \leanid{StagedContraction}: \leanid{edits_through_zero}, \leanid{two_contractions}, \leanid{no_single_contraction_to_core}, \leanid{signed_cores}, \leanid{factorable_by_signed_core}, \leanid{frozen_rejects}.
& \lc Sequences through zero at every tolerance; the starting support needs two peeling steps to reach its core;
the signed cores and the verdict; keeping the source equations for both
reductions leaves no common signed support.
& \lc The figure itself.\\
\leanrow
\lc Lemma~\ref{P-lem:limit}
& \lc \leanid{LimitContraction.lean}: \leanid{contracts_of_limit}.
\leanid{Components.lean}: \leanid{signedCore_of_chart_limit}, \leanid{retained_sign_of_limit} (private).
& \lc As stated.
& \lc None.\\
\leanrow
\lc Lemma~\ref{P-lem:interpolation}
& \lc \leanid{ChartPath.lean}: \leanid{fineConnected_same_chart}.
\leanid{PrimeInterpolation.lean}: \leanid{balanced_of_endpoints}.
\leanid{Mesh.lean}: \leanid{fineConnected_of_segment}.
& \lc For feasible rational endpoints of the same signed support,
\leanid{fineConnected_same_chart} concludes lawful \leanid{FineConnected}. Its
construction uses points of that signed chart.
& \lc The explicit restriction of every chain point to the common signed
chart is part of the written statement.\\
\leanrow
\lc Lemma~\ref{P-lem:compact-chain}
& \lc \leanid{CommitLateComponents.lean}: \leanid{fineMesh_iff_component} (private).
& \lc As stated, for a compact space and a map with dense range.
& \lc None.\\
\leanrow
\lc Theorem~\ref{P-thm:core}
& \lc \leanid{Classification.lean}: \leanid{fineConnected_iff_signedCore}.
\leanid{Components.lean}: \leanid{fineConnected_signedCore}, \leanid{realChart_closure_compact}.
& \lc As stated, with boundedness of the rational feasible locus as hypothesis.
& \lc None.\\
\leanrow
\lc Remark~\ref{P-rem:closure}
& \lc None.
& \lc ---
& \lc The whole remark.\\
\leanrow
\lc Corollary~\ref{P-cor:real}
& \lc \leanid{RealComponents.lean}: \leanid{real_joinedIn_iff_signedCore}, \leanid{real_connectedComponent_iff_signedCore}, \leanid{fineConnected_iff_real_joinedIn}.
& \lc In full.
& \lc None.\\
\bottomrule
\end{tabular}
\end{table}

\begin{table}[tbp]
\scriptsize\setlength{\tabcolsep}{3pt}
\caption{Lean correspondence for Section~\ref{P-sec:tractable}.}\label{tab:lean-tractable}
\begin{tabular}{@{}p{0.14\textwidth}p{0.33\textwidth}p{0.27\textwidth}p{0.20\textwidth}@{}}
\leanhead
\lc Lemma~\ref{P-lem:separate}
& \lc \leanid{FixedProfile.lean}: \leanid{FixedProfile.legal}.
\leanid{Capabilities.lean}, \leanid{Protocol}: \leanid{fiber}.
\leanid{Experiments.lean}, \leanid{NativeExperiments}: \leanid{protocol}, \leanid{physical_support}.
& \lc Lean represents the separate legality conditions, compiles the tests
and proves the table-support identity.
& \lc The combined formulation of Lemma~\ref{P-lem:separate} is written.\\
\leanrow
\lc Lemma~\ref{P-lem:transfer}
& \lc (a) \leanid{NativeProbes.lean}: \leanid{native_spectral_modulus}.
\leanid{FixedProfile.lean}: \leanid{spectralModulus_for_physicalStep}.
(c) \leanid{Rates.lean}: \leanid{rate_fineConnected}.
(d) \leanid{Register.lean}: \leanid{register_fineConnected}.
One-block estimates: \leanid{FixedProfile.lean}, \leanid{ParityBasis}: \leanid{kernel_sub_le}.
\leanid{FixedProfile.lean}: \leanid{finePhysical_of_signedCore}.
& \lc Parts (a), (c) and (d), and the one-block table estimates used in the
fixed-profile lift. For (c), Lean interpolates in inverse coordinates under
the bound $\sum_m|v_m|\le1$, not by normalized geometric interpolation.
& \lc The joint uniform-continuity statement in (b) is written.\\
\leanrow
\lc Remark~\ref{P-rem:restrictions}
& \lc None.
& \lc ---
& \lc The kernel example.\\
\leanrow
\lc Theorem~\ref{P-thm:certificates}
& \lc \leanid{CoreCertificate.lean}, \leanid{CoreCertificate}: \leanid{checkCore_sound}.
\leanid{CertificateComplete.lean}, \leanid{CoreCertificate}: \leanid{checkCore_complete_short}, \leanid{native_certificate_exists}.
\leanid{NativeCertificate.lean}, \leanid{CoreCertificate}: \leanid{nativeCheck_sound}, \leanid{nativeCheck_complete}.
\leanid{ComplexityTheorem.lean}, \leanid{EventExperiments}: \leanid{singleton_certificate}.
& \lc Soundness of supplied peeling and Gordan certificates, their existence,
and the stage count (at most one stage per removed coordinate). The
checker reads the source's compiled tests and both endpoints.
& \lc (a) the LP search and its time; (b) the multiplier and binary-length
bounds; (c) the decision time, certificate size and checking time.\\
\leanrow
\lc Proposition~\ref{P-prop:long-sequences}
& \lc None.
& \lc ---
& \lc The whole proposition.\\
\leanrow
\lc Theorem~\ref{P-thm:p-hard}
& \lc \leanid{MonotoneCircuitMarkov.lean}, \leanid{MonotoneCircuitMarkov.Circuit}: \leanid{native_fine_iff_value}, \leanid{endpoints_distinct}.
\leanid{MonotoneCircuitEncoding.lean}, \leanid{MonotoneCircuitMarkov.Encoding}: \leanid{checked_reduction}, \leanid{stream_length}, \leanid{fields_fit_bits}, \leanid{cursor_fits_bits}.
& \lc The circuit equivalence, with feasible endpoints at one common profile;
its transfer to the common word-test language; the indexed stream's record
count and its field and cursor-width bounds.
& \lc Goldschlager's theorem; full record decoding and the serialization
(Appendix~\ref{P-app:serialization}); the tape-level log-space
implementation.\\
\bottomrule
\end{tabular}
\end{table}

\begin{table}[tbp]
\scriptsize\setlength{\tabcolsep}{3pt}
\caption{Lean correspondence for Section~\ref{P-sec:tractable}, continued.}\label{tab:lean-tractable-cont}
\begin{tabular}{@{}p{0.14\textwidth}p{0.33\textwidth}p{0.27\textwidth}p{0.20\textwidth}@{}}
\leanhead
\lc Lemma~\ref{P-lem:covers}
& \lc \leanid{Effects.lean}: \leanid{Protocol.capKind_of_cover}.
& \lc For the underlying protocol, an attained cover changes exactly one
profile coordinate of the indicated kinds.
& \lc The classification for widened attained covers, the one-new-value-test
clause, and the admissible-intermediate-set criterion for word tests.\\
\leanrow
\lc Theorem~\ref{P-thm:directed}
& \lc \leanid{Experiments.lean}, \leanid{NativeExperiments}: \leanid{exact_six_iff}.
\leanid{ExactSix.lean}, \leanid{Protocol}: \leanid{factorable_iff}, \leanid{Chain.profile_le}.
\leanid{NativeCertificate.lean}, \leanid{CoreCertificate}: \leanid{ProfileTable.leCheck_correct}, \leanid{nativeCheck_sound}, \leanid{nativeCheck_complete}.
& \lc The directed signed-core equivalence, with both cores from the source's
constraints; the order comparison in the certificate check.
& \lc P-completeness, certificate size and checking time; the PSPACE-hardness
clause, from Theorem~\ref{P-thm:boundary}(ii) and
Remark~\ref{P-rem:inputs}(ii).\\
\leanrow
\lc The sign example (Sections~\ref{P-sec:directed} and~\ref{P-sec:acquisition})
& \lc \leanid{CommitLateExample.lean}, \leanid{CommitLateExample}: \leanid{early_acquisition}, \leanid{late_acquisition}, \leanid{late_factorable}, \leanid{early_chain_iff}, \leanid{early_to_source_no_chain}.
& \lc Fine reachability from $A$ to $B$; from $A^+$ a sequence to $B$ exists
exactly when $\varepsilon>1/8$; no sequence from $A^+$ back to $A$ at any
tolerance.
& \lc None.\\
\leanrow
\lc The four-label obstruction (Section~\ref{P-sec:directed})
& \lc \leanid{Example.lean}, \leanid{Example}: \leanid{obstructionTolerance_value}, \leanid{obstruction_at_native_tolerance}.
& \lc No sequence from the positive-$y$ endpoint to the negative-$y$
endpoint at tolerance $1/272$. The Lean catalogue also contains one
additional word test and one additional calibration, neither held at the
endpoints.
& \lc None.\\
\leanrow
\lc Remark~\ref{P-rem:readback}
& \lc \leanid{Readouts.lean}, \leanid{NativeExperiments}: \leanid{readouts_eq_iff}, \leanid{elementary_iff_step}, \leanid{elementary_iff_unique_effect}.
\leanid{QueryTables.lean}, \leanid{NativeExperiments}: \leanid{tables_eq_iff}, \leanid{table_chain_iff}, \leanid{table_exact_six}, \leanid{same_tests_same_effect}.
& \lc Within one fixed instrument, the readout table determines the
configuration, the correspondence preserves and reflects steps and
sequences, and the directed criterion is recoverable from the tables. Two
presentations into that instrument with matching tables have identical
elementary effect predicates.
& \lc The comparison of $A^+$ and $B$ in the sign example. Relabeling
across instruments under the stated attained-order, named-query,
final-level and projected-sum hypotheses.\\
\bottomrule
\end{tabular}
\end{table}

\begin{table}[tbp]
\scriptsize\setlength{\tabcolsep}{3pt}
\caption{Lean correspondence for Section~\ref{P-sec:hard}.}\label{tab:lean-hard}
\begin{tabular}{@{}p{0.14\textwidth}p{0.33\textwidth}p{0.27\textwidth}p{0.20\textwidth}@{}}
\leanhead
\lc Lemma~\ref{P-lem:engine}
& \lc \leanid{SingleAdditiveLicence.lean}, \leanid{SingleAdditiveLicence}: \leanid{joint_words_iff}, \leanid{floor}, \leanid{nclInstrument}.
& \lc The general measured comparison gives the mass and probe equations. In
the constructed model, common-profile feasibility forces the displayed
rates and loads.
& \lc The load-test iff for arbitrary shared-probe models with the stated
floor and calibrations.\\
\leanrow
\lc Lemma~\ref{P-lem:feasible}
& \lc \leanid{SingleAdditiveLicence.lean}, \leanid{SingleAdditiveLicence}: \leanid{measured_forces_rates} (private), \leanid{decode_lawful} (private), \leanid{encode_legal} (private), \leanid{encode_holds} (private).
\leanid{ConstraintLogic.lean}, \leanid{ConstraintLogic.Network}: \leanid{legal_iff_scaled}.
& \lc At the common profile the decoder obtains a lawful graph point. Under
the weight and capacity bounds the encoder gives a feasible configuration
from a lawful graph point, with clipped summands.
& \lc The iff for an arbitrary supplied parameter vector with all the
remaining model constraints retained is written.\\
\leanrow
\lc Lemma~\ref{P-lem:decode}
& \lc \leanid{ConstraintLogic.lean}, \leanid{ConstraintLogic.Graph}: \leanid{common_orientation}, \leanid{flips_of_represents}, \leanid{legal_of_represents}.
\leanid{SingleAdditiveLicence.lean}, \leanid{SingleAdditiveLicence}: \leanid{decode_chain} (private).
& \lc Common represented orientations at small summand steps, legality of
mixtures, and decoding of native chains at the cutoff.
& \lc The sequence conclusion under just the per-edge summand bound on each consecutive pair, without assuming the native cutoff.\\
\leanrow
\lc Lemma~\ref{P-lem:lift}
& \lc \leanid{ConstraintLogic.lean}, \leanid{ConstraintLogic.Graph}: \leanid{flip_segments_lawful}, \leanid{chain_of_flips}.
\leanid{SingleAdditiveLicence.lean}, \leanid{SingleAdditiveLicence}: \leanid{encode_chain} (private).
& \lc Feasibility of the three rational segments, their subdivision at
arbitrary positive tolerance, and transfer to native sequences.
& \lc None.\\
\leanrow
\lc Lemma~\ref{P-lem:cutoff}
& \lc \leanid{SingleAdditiveLicence.lean}, \leanid{SingleAdditiveLicence}: \leanid{ncl_chain_iff}, \leanid{ncl_fine_iff}, \leanid{tolerance}.
& \lc The equivalence at the supplied cutoff $\varepsilon_0$ and for fine
reachability.
& \lc Every smaller tolerance, by combining the two with monotonicity.\\
\leanrow
\lc The construction (Section~\ref{P-sec:construction}, Appendix~\ref{P-app:ncl-size})
& \lc \leanid{SingleAdditiveLicence.lean}, \leanid{SingleAdditiveLicence}: \leanid{ncl_exactly_one}, \leanid{ncl_licence_duration}, \leanid{ncl_base_size}, \leanid{ncl_endpoints_distinct}.
\leanid{BinaryParity.lean}, \leanid{BinaryParity}: \leanid{size_bounds}.
\leanid{ConstraintLogic.lean}, \leanid{ConstraintLogic.Network}: \leanid{flips_iff_scaled}.
& \lc Exactly one word-against-pair test and its length; the base-size bounds;
distinct constructed endpoints for distinct input orientations; the transfer from weighted networks to scaled charges.
& \lc Hearn and Demaine's Corollary~6; the binary enumeration of labels; the
running time and binary-length bounds.\\
\bottomrule
\end{tabular}
\end{table}

\begin{table}[tbp]
\scriptsize\setlength{\tabcolsep}{3pt}
\caption{Lean correspondence for Section~\ref{P-sec:hard}, continued.}\label{tab:lean-hard-cont}
\begin{tabular}{@{}p{0.14\textwidth}p{0.33\textwidth}p{0.27\textwidth}p{0.20\textwidth}@{}}
\leanhead
\lc Proposition~\ref{P-prop:length-one}
& \lc \leanid{SingleAdditiveLicence.lean}, \leanid{SingleAdditiveLicence}: the \leanid{short} parameter of \leanid{nclInstrument} and of each theorem above.
& \lc Each length separately.
& \lc The cancellation identifying the two feasible sets.\\
\leanrow
\lc Remark~\ref{P-rem:control}
& \lc \leanid{SingleAdditiveLicence.lean}, \leanid{SingleAdditiveLicence}: \leanid{control_replacement}, \leanid{Control.endpoint_profile}, \leanid{Control.fine}, \leanid{Control.verdict_changes}.
& \lc The replacement keeps the endpoints and their profile; a yes verdict at
every tolerance; the change of verdict, assuming that no legal flips join
the orientations.
& \lc The $K_{3,3}$ example, which supplies that assumption.\\
\leanrow
\lc Corollary~\ref{P-cor:real-complexity}
& \lc \leanid{RealAlgebraicBoundary.lean}, \leanid{MonotoneCircuitMarkov.Circuit}: \leanid{algebraic_real_path_iff_value}.
\leanid{RealAlgebraicBoundary.lean}, \leanid{AlgebraicConstraintLogic}: \leanid{real_path_iff_flips}, \leanid{flips_of_real_path}, \leanid{network_real_path_iff_flips}.
\leanid{RealComponents.lean}: \leanid{real_joinedIn_iff_signedCore}.
& \lc The real signed-core criterion without boundedness; the circuit
equivalence; both directions of the real NCL correspondence for the
constructed endpoints.
& \lc The LP resource bounds, input encodings and machine-level reductions;
the upper bound (Appendix~\ref{P-app:hrs}).\\
\leanrow
\lc Theorem~\ref{P-thm:rational-roots}
& \lc \leanid{ResponseBudgets.lean}, \leanid{ResponseBudgets}: \leanid{assign_lawful}, \leanid{decode_solution}, \leanid{decode_lift}, \leanid{lift_chain}, \leanid{decode_chain}, \leanid{fine_iff_box}, \leanid{no_native_chain_of_selector}.
\leanid{CalibratedResponses.lean}, \leanid{CalibratedResponses}: \leanid{rates_eq}, \leanid{lawful_state}, \leanid{build_decode}.
& \lc For supplied response rows, a parity basis and feasible box endpoints,
Lean verifies the row compiler and the transfers in both directions. The
fixed-tolerance obstruction assumes that every feasible box solution has
selector zero or one and that the endpoints have opposite selector values.
& \lc The selector dichotomy; the gate compilation of sparse polynomials; the
compiler's binary-length bounds and running time; the residual-to-row
assembly and the three-segment construction.\\
\leanrow
\lc Corollary~\ref{P-cor:quartic}
& \lc \leanid{QuarticSelector.lean}, \leanid{QuarticSelector}: \leanid{homogeneous_scale}, \leanid{interior_hasRoot}, \leanid{no_root_selector}, \leanid{root_has_three_segments}, \leanid{RowRealization.native_boundary}, \leanid{RowRealization.nativeEndpoints_distinct}.
\leanid{ResponseBudgets.lean}, \leanid{ResponseBudgets}: \leanid{square_rows}.
& \lc The homogeneous identity, root extraction and three segments; given a
\leanid{RowRealization} (lift, readback and slope proofs), both verdicts at
the exact tolerance, exactly one word-against-pair test and distinct
endpoints.
& \lc The fixed-layout row assembly, its counts and binary-length bounds; the conditional
consequence.\\
\bottomrule
\end{tabular}
\end{table}

\begin{table}[tbp]
\scriptsize\setlength{\tabcolsep}{3pt}
\caption{Lean correspondence for Section~\ref{P-sec:timing}.}\label{tab:lean-timing}
\begin{tabular}{@{}p{0.14\textwidth}p{0.33\textwidth}p{0.27\textwidth}p{0.20\textwidth}@{}}
\leanhead
\lc Theorem~\ref{P-thm:acquisition}
& \lc \leanid{CommitLate.lean}, \leanid{NativeExperiments}: \leanid{commitLate_iff}.
& \lc The equivalence, for equal parameter vectors and ordered profiles.
& \lc The polynomial-time claim.\\
\leanrow
\lc Theorem~\ref{P-thm:projected-components}
& \lc (a) \leanid{CommitLate.lean}, \leanid{EventExperiments}: \leanid{carry}, \leanid{atSource}, \leanid{projected_near}, \leanid{carry_near}, \leanid{fixedProfile_registerChain}, \leanid{chain_of_sameNumbers}, \leanid{factorable_iff_atSource}.
(b) \leanid{CommitLateComponents.lean}, \leanid{EventExperiments}: \leanid{ProjectedClosure}, \leanid{fixedProfile_iff_projectedComponent}, \leanid{factorable_iff_projectedComponent}.
\leanid{EventExperiments.lean}, \leanid{EventExperiments}: \leanid{Chain.profile_le}.
(c) \leanid{CommitLateComponents.lean}, \leanid{EventExperiments}: \leanid{projectedClosure_mono}.
(d) \leanid{CommitLateComponents.lean}, \leanid{EventExperiments}: \leanid{commitLate_iff_projectedComponents}.
& \lc (a)--(d). The Lean form of (d) allows different parameter vectors at $A$
and $A^+$ and assumes no cover.
& \lc None.\\
\leanrow
\lc Proposition~\ref{P-prop:source-schedule}
& \lc \leanid{CommitLate.lean}, \leanid{EventExperiments}: \leanid{commitAtTargetNumbers}, \leanid{commitAtTarget}, \leanid{commitAtTargetNumbersAlong}.
& \lc Given fine reachability from $A$ to $B$, Lean constructs stages at the
prescribed ordered profiles with the target's label probabilities and
eigenvalues. Each stage is finely reachable from $A$ and can finely reach
$B$, and every ordered pair of stages has the corresponding
fine-reachability relation. A stage at a profile with the target's support
can use the target's entire parameter vector.
& \lc Concatenation into one sequence at a specified tolerance.\\
\bottomrule
\end{tabular}
\end{table}

\begin{table}[tbp]
\scriptsize\setlength{\tabcolsep}{3pt}
\caption{Lean correspondence for Section~\ref{P-sec:timing}, continued, including the results stated in Appendix~\ref{P-app:core-order-proofs}.}\label{tab:lean-timing-cont}
\begin{tabular}{@{}p{0.14\textwidth}p{0.33\textwidth}p{0.27\textwidth}p{0.20\textwidth}@{}}
\leanhead
\lc Theorem~\ref{P-thm:no-calibration}
& \lc \leanid{NativeCoreOrder.lean}: \leanid{calibration_mem_core}.
\leanid{NativeCoreOrder.lean}, \leanid{NativeExperiments}: \leanid{active_spectral_rows_balanced}, \leanid{spectral_core_empty_iff}, \leanid{signature_eq_zero_iff}, \leanid{factorable_iff_profile_le_of_no_nonzero_calibration}, \leanid{realise_zero_of_no_nonzero_calibration}.
\leanid{CoreOrderExamples.lean}, \leanid{CoreOrderExamples}: \leanid{product_not_fineConnected}, \leanid{Native.zero_calibration_formal_support}.
& \lc As stated, including the configuration with zero eigenvalues; the two
examples after it.
& \lc None.\\
\leanrow
\lc Proposition~\ref{P-prop:empty-core}
& \lc \leanid{CoreOrder.lean}: \leanid{core_eq_empty_iff}, \leanid{core_eq_empty_of_positive_damping}, \leanid{core_eq_empty_of_balanced_active}, \leanid{fineConnected_of_empty_cores}.
& \lc (a)--(c).
& \lc None.\\
\leanrow
\lc Corollary~\ref{P-cor:refine}
& \lc \leanid{CoreOrder.lean}: \leanid{Contracts.of_rows_subset}, \leanid{Reduces.of_rows_subset}, \leanid{core_mono_rows}, \leanid{core_absorb_rows}, \leanid{signedCore_absorb_rows}, \leanid{signedCore_eq_of_rows_subset}, \leanid{FineConnected.of_tests_subset}.
\leanid{NativeCoreOrder.lean}: \leanid{Protocol.spectralRows_mono}, \leanid{NativeExperiments.factorable_of_higher_source}.
\leanid{CoreOrderExamples.lean}, \leanid{CoreOrderExamples.Native}: \leanid{strict_core_growth}, \leanid{strict_reachability_split}.
& \lc (a) and (b); the example of strict growth.
& \lc None.\\
\leanrow
\lc Proposition~\ref{P-prop:interior}
& \lc \leanid{CoreOrder.lean}: \leanid{coreInterior}, \leanid{removedClosure}, \leanid{removedClosure_isClosed_iff}, \leanid{core_greatest_irreducible}.
\leanid{NativeCoreOrder.lean}, \leanid{NativeExperiments}: \leanid{realise_core}, \leanid{attainedCoreInterior}, \leanid{core_greatest_attained_irreducible}.
\leanid{CoreOrderExamples.lean}, \leanid{CoreOrderExamples.Native}: \leanid{core_not_extensive_on_attained_supports}.
& \lc As stated; the example after it.
& \lc None.\\
\bottomrule
\end{tabular}
\end{table}

\begin{table}[tbp]
\scriptsize\setlength{\tabcolsep}{3pt}
\caption{Lean correspondence for Section~\ref{P-sec:exactness}, including the
results stated in Appendix~\ref{P-app:zero-probes}.}\label{tab:lean-exactness}
\begin{tabular}{@{}p{0.14\textwidth}p{0.33\textwidth}p{0.27\textwidth}p{0.20\textwidth}@{}}
\leanhead
\lc Proposition~\ref{P-prop:one-probe}
& \lc \leanid{FermatProbe.lean}, \leanid{FermatProbe}: \leanid{corner_lawful}, \leanid{endpoint}, \leanid{no_chain}, \leanid{not_factorable}, \leanid{exactly_one}, \leanid{realArc}, \leanid{continuous_realArc}, \leanid{realArc_endpoints}, \leanid{realArc_quartic}, \leanid{quartic} (private), \leanid{axis} (private), \leanid{marked_diag} (private).
& \lc The feasible endpoints, no sequence at tolerance $1/72$, and a continuous
real arc satisfying the equations; uses Mathlib's Fermat theorem for
exponent four, \leanid{fermatLastTheoremFour}.
& \lc The real arc as a path of real-valued configurations (Lean's
configurations are rational).\\
\leanrow
\lc Proposition~\ref{P-prop:zero-probe-fermat}
& \lc None.
& \lc ---
& \lc The whole proposition, with Darmon and Granville's theorem.\\
\leanrow
\lc Proposition~\ref{P-prop:separation}
& \lc \leanid{CoreGap.lean}, \leanid{CoreCertificate}: \leanid{dual_log_identity}, \leanid{dual_survivor_lower}.
\leanid{CoreGap.lean}, \leanid{Protocol}: \leanid{signature_eq_of_bounded_mesh}.
& \lc The separation lemmas under supplied bounds: the lower bound assumes an
accepted Gordan certificate and the radius, coordinate and coefficient
bounds; the model statement assumes a uniform bound.
& \lc The radius, the determinant estimates and $D$; part~(c); the
binary-exponent variant.\\
\leanrow
\lc Proposition~\ref{P-prop:zero-probes}
& \lc \leanid{NoProbeRatios.lean}, \leanid{NoProbeRatios}: \leanid{rank_one_no_three}.
& \lc The strict-convexity step of the rank-one case.
& \lc Everything else, including the polynomial-time claim.\\
\leanrow
\lc Proposition~\ref{P-prop:zero-probe-closure}
& \lc None.
& \lc ---
& \lc The whole proposition.\\
\leanrow
\lc Proposition~\ref{P-prop:one-state}(a)
& \lc \leanid{BooleanEvents.lean}, \leanid{OneStateEvents.BooleanReduction}: \leanid{native_fine_iff_flips}, \leanid{native_chain_iff_flips_of_le}, \leanid{endpoints_distinct}.
\leanid{ClauseReconfiguration.lean}, \leanid{ClauseReconfiguration}: \leanid{fine_iff_flips}.
\leanid{EventCompiler.lean}, \leanid{OneStateEvents.Equations}: \leanid{mark_card}, \leanid{tests_length}, \leanid{tests_ne_nil}, \leanid{encode_lawful}, \leanid{lawful_decode}, \leanid{fine_encode_iff}.
\leanid{WidenedOneState.lean}, \leanid{OneStateEvents}: \leanid{widened_chain_iff}.
\leanid{ComplexityTheorem.lean}, \leanid{OneStateEvents.BooleanReduction}: \leanid{widened_fine_iff_flips}, \leanid{widened_chain_iff_flips_of_le}, \leanid{widened_endpoints_distinct}.
& \lc The finite compiler correspondences for tautology-free formulas with
satisfying endpoints, at every tolerance at most $2\delta$ and for fine
reachability; the transfer to the model with named word tests. Lean also
checks the compiler's exact label and test counts for supplied
expression-node data.
& \lc The imported PSPACE-hardness and its preprocessing; the clause-template
sizes and total binary-output bound, encodings and running time.\\
\leanrow
\lc Proposition~\ref{P-prop:one-state}(b)
& \lc \leanid{NativeTractable.lean}, \leanid{OneStateEvents}: \leanid{one_step_event_fine}, \leanid{linearTests}, \leanid{singleton_event_fine}, \leanid{SingletonTest}, \leanid{singletonTests}.
& \lc For external predicates consisting of equalities of sums over label
sets, Lean proves fine connectivity under a positive floor. It separately
proves fine connectivity for lists of equal-length monomial-ratio tests
with positive ratios.
& \lc The general case of all held equations linear in the rates, including
mixed duration-one word tests and probability-ratio tests in the paper's
named model.\\
\leanrow
\lc Proposition~\ref{P-prop:one-state}(c)
& \lc \leanid{ArithmeticEvents.lean}, \leanid{OneStateEvents.ArithmeticReduction}: \leanid{native_fine_iff_solvable}, \leanid{native_chain_iff_solvable_of_le}, \leanid{endpoints_distinct}.
\leanid{PolynomialEvents.lean}, \leanid{OneStateEvents.ArithmeticReduction.Term}: \leanid{solvable_iff_zero}, \leanid{native_fine_iff_rational_zero}, \leanid{native_chain_iff_rational_zero_of_le}.
& \lc The reduction for supplied gate lists and arithmetic terms, on auxiliary
one-state configurations defined by external test predicates.
& \lc The sparse-input and test-list construction, template bounds, running
time; path-connectedness of the real feasible set.\\
\bottomrule
\end{tabular}
\end{table}
% lint: lean-table end

\end{document}
